% !TEX root = ../main.tex
% Denne fil kompileres via main.tex.
\documentclass[../main.tex]{subfiles}
\usepackage{amsmath,amssymb}
\begin{document}

\section{Numerisk ustabilitet}

$C$ er korrelationsmatricen

Numerisk ustabilitet har stor betydning indenfor moderne portefølje teori. Da det forudsætter at vores traditionelle optimeringsmetoder ikke er tilstrækkelige for at kunne sikre et konsistent og rigtigt resultat.

I denne opgave ser behandler vi numerisk ustabilitet i populationens kovarians- og korrelationsmatrice. Problemet er at disse er ukendte og vi forsøger at estimere disse gennem stikprøvekovarians- og stikprøvekorrelationsmatricen. I estimationsprocessen opstår der støj der bliver inkluderet i stikprøven. Der er 2 typer for støj der påvirker vores estimation at de brugte matricer. Det første element er støj fra stikprøveusikkerhed. Den kan minimeres ved at formindske $N/T$ ratioen, hvor $T$ er antallet af observationer og $T$ er antallet af aktiver. Den anden fejlkilde er støj fra datastrukturen. Her påvirker krydskorrelerede aktiver rangen af kovariansmatricen, og dermed condition number. Det medfører signalustabilitet der ikke kan vi formindskes ved en formindskning af $N/T$ ratioen.


\subsection{Random Matrix Theory}

For at kunne specificere hvordan numerisk ustabilitet fremkommer af støj, introducerer vi Random Matrix Theory.

Vi vælger at betragte aktiemarked med den bagvedliggende datagenerende proces som gaussiske observationer med middelværdi$=0$ og varians$=\sigma^2$. Hvor $0<q=N/T\le1$. Samt estimerer vi stikprøvekovariansmatricen $\Sigma=T^{-1}X'X$.

Marchenko-Pastur theorem viser for $N\to\infty$ og $T\to\infty$. Samt antagelser om I.I.D at fordelingen for egenværdierne for stikprøvekovariansmatricen konvergerer mod Marchenko-Pastur fordelingen:

\begin{equation}
    \rho_mp[\lambda] =
    \begin{cases}
        \dfrac{T}{N}\dfrac{\sqrt{(\lambda_{+}-\lambda)(\lambda-\lambda_{-})}}{2\pi\lambda\sigma^{2}} & \text{hvis } \lambda \in [\lambda_{-},\lambda_{+}] \\
        0 & \text{hvis } \lambda \notin [\lambda_{-},\lambda_{+}]
    \end{cases}
    \label{eq:mp_density}
\end{equation}

Hvor den forventede værdi for minimum og maximum af egenværdierne er:
\begin{equation}
    \lambda_{\pm} = \sigma^{2} \left( 1 \pm \sqrt{q} \right)^{2}
    \label{eq:mp_bounds}
\end{equation}

Under nulhypotesen netop $C=I$, ligger egenværdierne indenfor intervallet $[\lambda_-,\lambda_+]$ og opstår derfor alene af estimationsstøj. I en empirisk stikprøvekovariansmatrice er ikke alle egensystemer fuldkommen tilfældige. Således er det kun egenværdierne udenfor intervallet for nulhypotesen der afgiver et signifikant statisk signal. For finansielle dataset er signal til støj ratioen kendt for at være lav som et resultat af arbitrære krafter.\cite{lopez_de_prado_machine_2020}
\begin{figure}
    \centering
    \includegraphics[width=0.5\linewidth]{Teori/MP fordeling med data.jpg}
    \caption{Marchenko-Pastur fordelingen med empirisk data}
    \label{fig:placeholder}
\end{figure}



Ser vi kort på matrice invertering, kan vi bruge samme udgangspunkt i MP fordelingen til at forstå hvordan støj forværrer estimering til optimeringsmetode. 
 Ved asymptotik og $\hat{\Sigma}=I$ gælder det at stikprøvekovariansmatricens  egenværdier vil  konvergere mod MP-fordelingen. Dette har den direkte implikation for den numeriske stabilitet af kovariansestimatet. Når q-rationen ikke er ubetydelig lille vil $\lambda_-\to0$ når $q\to1$, som fremkommer af $\lambda_-=\sigma^2(1-\sqrt{q})^2$. Dette medfører at condition number for stikprøvematricen der er defineret som
\begin{equation}\kappa(\hat{\Sigma})=\lambda_{\max}/\lambda_{\min}\approx\lambda_+/\lambda_-=((1+\sqrt{q})^2)/((1-\sqrt{q})^2)\end{equation}
Vokser eksplosivt når q går mod 1. Et højt Konditionstal indebærer at matricen er ill-conditioned. Altså små ændringer i data som i vores tilfælde kunne være støj, giver store udsving i kovariansmatricens inverse. 

Dette kan direkte observeres ved spektraldekompositionen. $\hat{\Sigma}=Q\Lambda Q'$ og hvor $\hat{\Sigma}^{-1}=Q\Lambda^{-1}Q'$. Hvor den inverse matrice bruges til klassiske porteføljeoptimeringsmetoder. Udfordringen for disse metoder er at da $\Lambda^{-1}=\mathrm{diag}(1/\lambda_1,\ldots,1/\lambda_N)$ hvor $1/\lambda_i$ bliver meget stor da egenværdierne er tæt på 0. Dermed forstærkes estimationsstøjen der resulterer i numerisk ustabilitet.







\subsection{Numerisk ustabilitet som følge af signal}


Selve signalet kan også medvirke til at der opstår numerisk ustabilitet for korrelationsmatricen. Vi tager igen udgangspunkt  i C der er vores korrelationsmatrice der er en 2x2 matrice. 
$\begin{equation}
    C = \begin{bmatrix}
        1 & \rho \\
        \rho & 1
    \end{bmatrix}
    \label{eq:placeholder_label}
\end{equation}$

C matricen kan blive diagonaliseret som $CW=W\Lambda$ hvor

\begin{equation}
    \Lambda = \begin{bmatrix} 1 + \rho \\ 1 - \rho \end{bmatrix} ;\;
    W = \begin{bmatrix}
        \dfrac{1}{\sqrt{2}} & \dfrac{1}{\sqrt{2}} \\
        \dfrac{1}{\sqrt{2}} & -\dfrac{1}{\sqrt{2}}
    \end{bmatrix}
    \label{eq:placeholder_label}
\end{equation}

Den inverse af C er $C^(-1)=W\Lambda^-1W'= \frac{1}{|C|} 
    \begin{bmatrix}
        1 & -\rho \\
        -\rho & 1
    \end{bmatrix}$
Hvor $|C|$ er determinanten af C. 

I tilfælden hvor $\rho=0$, hvor aktiverne er ukorrelererede så er  $C=I$. Dette svarer til at egenværdier ligger i en horisontal linje hvor egenvektorerne $\lambda_1=\lambda_2=1$.Men i tilfældet hvor aktiverne ikke er ukorrelerede bliver egenværdierne $\lambda_1=1+\rho,\lambda_2=1-\rho$
I realiteten vides der at i finansielle markeder opstår der klyngekorrelationer, der svarer til at nogle aktiver følger hinanden mere end andre aktiver. Det kunne for eksempel være aktiver der ligger i samme sektor, som tech- eller medicinalbranchen. Vi undersøger hvad der sker når der i eksemplet bliver dannet en klyngekorrelation mellem de 2 aktiver. Når $\rho\rightarrow1$, har det effekten, $\lambda_1\rightarrow1, \lambda_2\rightarrow0$. Dermed bliver determinanten $Det(C)=(1+\rho)(1-\rho)=1-\rho^2$ og i tilfældet hvor aktiverne går i mod fuld korrelerede. $\rho\rightarrow1, Det(C)\rightarrow0$. Dermed nærmer matricen sig singularitet. Der følger at inversen igen er yderst følsom overfor små ændringer i korrelationsmatricen. 


Hvis vi fortsætter med brugen af eksemplet i højere dimensioner, kan vi bedre observere hvordan klyngerne påvirker resten af matricen. Når der opstår en klynge har det af betydning at  
\begin{equation}
    \label{eq:placeholder_label}
    \operatorname{tr}(C) = \sum_{i=1}^{N} \lambda_i=N
\end{equation}
Dermed hvis en klynge bestående af K akiver bliver påvirket af den samme fælles egenvektor hvor den tilhørende egenværdi forklarer en større del af den samlede varians. Men da sporet af C er fast, betyder det at som konsekvens at K-1 egenværdier i klyngen vil blive dæmpet, da signalet koncentrerer varians i få retninger. Konsekvensen bliver at jo større den indre klyngekorrelation er desto højere bliver konditionstallen $\kappa(C)$. Dermed er denne form for numerisk ustabilitet urelaterbar med numerisk ustabilitet af støj, da vi ikke kan formindske effekten ved at øge antallet af datapunkter.\cite{lopez_de_prado_robust_2019}













\end{document}
